\section{정리의 증명}
\label{sec:proofs}

\subsection{정리~\ref{thm:optimal_hamiltonian}의 증명}
\label{sec:proof_optimal}

\begin{proof}
위치 $\mathrm{pos}$에서 Post-RoPE 키 벡터 $k_{\mathrm{post}} = R_{\mathrm{RoPE}} k_{\mathrm{pre}}$의 공분산은
\[
\Sigma(\mathrm{pos}) = R_{\mathrm{RoPE}} \Sigma_0 R_{\mathrm{RoPE}}^\top
\]
이다. 여기서 $R_{\mathrm{RoPE}} = e^{\mathrm{pos} \cdot A_{\mathrm{RoPE}}}$이고 $\Sigma_0 = \E[k_{\mathrm{pre}} k_{\mathrm{pre}}^\top]$는 Pre-RoPE 공분산이다.

\textbf{Step 1.} Class~C 내에서 MSE를 최소화하는 회전은 공분산의 고유벡터 행렬이다 (정리~\ref{thm:master_formula}에서 $G(d)$ 달성 조건). $\Sigma(\mathrm{pos})$의 고유분해는
\[
\Sigma(\mathrm{pos}) = (R_{\mathrm{RoPE}} V_0) \Lambda (R_{\mathrm{RoPE}} V_0)^\top
\]
이므로, 위치 $\mathrm{pos}$에서의 최적 회전은 $U^*(\mathrm{pos}) = R_{\mathrm{RoPE}}(\mathrm{pos}) \cdot V_0$이다.

\textbf{Step 2.} 실제 구현에서는 Post-RoPE 키에 $U^*(\mathrm{pos})^\top$를 곱한다:
\[
U^*(\mathrm{pos})^\top k_{\mathrm{post}} = V_0^\top R_{\mathrm{RoPE}}^\top \cdot R_{\mathrm{RoPE}} k_{\mathrm{pre}} = V_0^\top k_{\mathrm{pre}}.
\]
\textbf{RoPE가 상쇄된다.} 따라서 위치에 무관하게 $V_0$를 Pre-RoPE 키에 적용하면 모든 위치에서 동시에 최적이 된다.

\textbf{Step 3.} $U^* = R_{\mathrm{RoPE}} V_0 = e^{\mathrm{pos} \cdot A_{\mathrm{RoPE}}} V_0$는 $A_{\mathrm{PCA}}(\Sigma_0)$에 대한 수반 표현 $\mathrm{Ad}_{R_{\mathrm{RoPE}}}(A_{\mathrm{PCA}})$의 지수 사상이다. MSE 이득은 $R_{\mathrm{aniso}} = \mathrm{AM}(\lambda)/\mathrm{GM}(\lambda)$이다 (따름정리~\ref{cor:gain_over_turbo}).
\end{proof}

\subsection{따름정리~\ref{cor:post_rope_fails}의 해설}

Post-RoPE PCA는 $\Sigma_{\mathrm{mixed}} = \frac{1}{T}\sum_t \Sigma(t)$에 대해 고유분해한다. $\Sigma(t)$가 위치에 따라 달라지므로, $\Sigma_{\mathrm{mixed}}$의 고유벡터는 어떤 특정 위치의 최적 기저와도 일치하지 않는다. 특히 RoPE 주파수가 기하급수적으로 퍼져 있으므로, 고주파 성분의 혼합이 심하다. 실측에서 Post-RoPE PCA+WF의 2-bit MSE가 TurboQuant보다 높은 것(Qwen: 1.091 vs 0.763)은 이 현상을 반영한다.

\subsection{Water-Filling floor=2의 이론적 동기}

$B$-bit 균일 양자화기의 MSE 비율 $c(B) = D(B) / \sigma^2$ (가우시안 소스)는:

\begin{center}
\begin{tabular}{@{}cccc@{}}
\toprule
$B$ & 1-bit & 2-bit & 3-bit \\
\midrule
$c(B)$ & 0.3634 & 0.1175 & 0.03454 \\
\bottomrule
\end{tabular}
\end{center}

1-bit에서 $c(1) = 0.36$은 분산의 36\%만 제거하므로, 1-bit 채널에 할당된 비트의 ``효율''이 매우 낮다. 반면 2-bit는 88\% 제거로 합리적이다. Qwen2.5-7B 2-bit에서의 floor ablation:

\begin{center}
\begin{tabular}{@{}ccccc@{}}
\toprule
floor & 1 & 2 & 3 & Uniform \\
\midrule
PPL & 11.255 & 7.099 & 6.757 & 7.980 \\
\bottomrule
\end{tabular}
\end{center}

floor=1은 치명적이고 (uniform 대비 41\% 악화), floor$\geq$2는 모두 uniform보다 개선된다. floor=3이 floor=2보다 약간 좋지만, 총 비트 예산 제약으로 인해 고분산 채널에 할당 가능한 비트가 줄어들어 trade-off가 존재한다. 본 논문에서는 보수적으로 floor=2를 기본값으로 사용한다.
